\begin{figure}[p]\centering
\includegraphics[width=6.4in]{figures/faraday_instrument_leakage.pdf}
\caption{Instrument and diffuse leakage. I-only native-frequency Haslam+CMB harmonic predictions at sidereal phase 0 through the as-built, C4-symmetrized and diagonal-ZA responses. Every variant uses the as-built band's fixed zenith calibration and enclosure normalization, plus identical MEPA time averaging. Curves show actual native frequencies; fine-channel linear/PCHIP leakage and their common-fit residuals are model interpolants exported in the result record. Variant differences are a model spread, not a beam-error probability distribution. These diffuse-sky leakages need not vanish after zenith calibration.}
\end{figure}
\begin{figure}[p]\centering
\includegraphics[width=6.4in]{figures/paper_depth_depolarization.pdf}
\caption{Depth-profile benchmark. Conditional flat-plane depth-profile moments. Smooth \textbar{}E[P]\textbar{} squared and centered speckle power are divided by the same unrotated beam-weighted reference. The Gaussian profile is supported on f in [0,1]; fractions 0/0.6/1 retain the original screen limits. Monochromatic points use 10/30/50 MHz; lines guide the eye. Triangles mark values below the 1e-12 display floor, including exact zeros and underflow. The source/beam geometry and local screen parameters remain conditional. This independent benchmark includes no real four-port beam, channel averaging, leakage or receiver noise; those enter the subsequent instrument predictions.}
\end{figure}
\begin{figure}[p]\centering
\includegraphics[width=6.4in]{figures/faraday_redistribution_spin2.pdf}
\caption{Real-instrument redistribution. Conditional as-built loaded four-port predictions at the reference MEPA sidereal phase, using the two illustrative local Gaussian columns and the five shared-column emitting-depth profiles. The first two rows show Var(Q)+Var(U) for the phase-mean and centered terms, divided by the same unrotated-beam reference power. The bottom row is sqrt[Var(Q)+Var(U)] in K. Points use their own native 10/30/50 MHz beams; connecting segments guide the eye. Triangles mark zero/underflow values at the 1e-16 display floor. The intrinsic spin-2 RMS is 1 K at 30 MHz with assumed E/B spectra, beta=-2.8 and a 15 MHz turnover. This is an illustrative brightness law, especially below 20 MHz. Slab/Gaussian emission depths share one column rather than independent rotating layers. These are model second moments, without I-only leakage or measurement noise.}
\end{figure}
\begin{figure}[p]\centering
\includegraphics[width=6.4in]{figures/faraday_measured_spectra.pdf}
\caption{Simulated measured spectra. Simulated observations, not mission data: as-built calibrated Q/I and U/I at preselected sidereal phase 0, uniform emitting slab and sigma\_phi=0.02 rad/m$^2$. Zoom curves retain all folded images; open circles are three parent measurements. A constant and physical-frequency slope pass through each channel response before fitting. Residuals are subtracted in K, then divided by the same fixed channelized sky-I reference as the components. Estimated-I covariances, including I-Q/U correlations, are exported. The paired I-only control shares the noise draw. The last row shows the underlying smooth and speckle polarization draws; only noise changes between repeated visits. Intrinsic RMS is 1 K at 30 MHz; local structure, turnover, current noise and temperatures are assumed. Voltage/current references are 3 nV/sqrt(Hz), 0.1 pA/sqrt(Hz). Exposure spans 24 illustrative lunations at 50\% duty. The colored band shows marginal conditional sky-plus-noise RMS around the residual I-only model (dashed). Native interpolation limits fine-channel leakage interpretation.}
\end{figure}
\begin{figure}[p]\centering
\includegraphics[width=6.4in]{figures/faraday_channel_covariance.pdf}
\caption{Frequency correlations. Conditional slab-signal Q correlations before and after the common smooth fit, on all 192 nominal zoom frequencies, for the same model and time bin as the measured-spectra figure. Before-fit images show 1-r on a common log scale with floor 1e-12; after-fit images show signed r. The last row retains individual Q and U correlations to the named center-frequency bin; unlike frequency pairs are not averaged into a stationary estimate. The noise-control line includes channel folds, overlap and the same fit. Full Q-U cross covariance and both parent/zoom blocks are exported. Repeated-cycle differences and their radiometer prediction are shown in faraday\_repeatability.pdf. These conditional correlations are not a Faraday identification or an observed covariance estimate.}
\end{figure}
\begin{figure}[p]\centering
\includegraphics[width=6.4in]{figures/faraday_channel_covariance_amplitude.pdf}
\caption{Q/U covariance amplitudes. Conditional total polarized-signal Q and U covariances on all 192 nominal zoom frequencies, before and after the common fit, for the same reference slab and phase as the measured spectra. The upper four panels share one positive logarithmic scale; the lower four share one signed linear scale. Each is in K squared at the common loaded calibration plane. The before/after groups use different scales because the fit removes most of this coherent reference signal. These are ensemble signal covariances, not covariances estimated from a single spectrum. The separate noise control and full Q-U, parent/zoom and complex C/J blocks are exported; normalized correlations are shown in the companion.}
\end{figure}
\begin{figure}[p]\centering
\includegraphics[width=6.4in]{figures/faraday_separability.pdf}
\caption{Conditional separability. Conditional residual Q/U RMS divided by the same fitted radiometer RMS over all signed uniform mean columns from the adopted density/field/path pairs. Every plotted mean includes both assumed dispersions, Gaussian and global-variance-mixture statistics, all five depth profiles and nine receiver-noise combinations. Solid points are maximum speckle ratios; shading spans positive speckle ratios. The unrotated f=0 case has exactly zero speckle. Dashed open squares give maximum total (smooth plus speckle) ratios, including f=0. These are conditional model ranges, not confidence intervals. The horizontal dotted line is equal RMS, not a detection threshold. The full arrays retain both phases and 1/24-lunation availability; this figure uses phase 0 and 24 illustrative lunations. Intrinsic RMS is 5 K at 30 MHz, the assumed volume-averaged polarized emissivity times the full 5 pc path. The depth integral supplies depolarization; no second horizon cut is made. Alternative emissivity normalizations are tabulated separately and scale amplitudes linearly.}
\end{figure}
\begin{figure}[p]\centering
\includegraphics[width=6.4in]{figures/faraday_repeatability.pdf}
\caption{Repeated sidereal-phase measurements. Two simulated repeat measurements of the same sky and instrument at sidereal phase 0. Top: Q residuals, with parent points and zoom curves. Bottom: their difference divided by the predicted independent-repeat radiometer RMS, sqrt(2) times the individual RMS; gray is $\pm$1. The sky, smooth polarization, speckle and leakage are held fixed. Noise draws are independent; each cycle uses one eligible phase-bin exposure at 50\% duty, exported in seconds. Correlated power sampling uncertainties and independent mock observations are recorded, without counting channels as independent skies.}
\end{figure}
\begin{figure}[p]\centering
\includegraphics[width=6.4in]{figures/faraday_angular_and_fringe.pdf}
\caption{Angular and observed fringe spectra. Conditional unrotated, foreground and uniform-slab polarized angular spectra (top), and their Q fringe power through a common instrument and observing operator (bottom). All use the local 0.02 rad/m$^2$ dispersion illustration and 1 K intrinsic RMS at 30 MHz. The operator averages 256 sidereal bins, selects complete actual-site nightside bins and applies their normalized Fourier transform. The same mask and finite-bin average, with common harmonic cutoff ell=48, act on polarized covariance, deterministic Haslam I leakage and explicit loaded receiver/radiometer noise. Noise here uses a 25 kHz top-hat at 50\% duty per retained bin. The window mixes modes: low-ell input shape alone does not establish separability, and no angular-map sensitivity is inferred. Full complex fringe covariance and relation are exported. Angular power below the displayed twelve-decade range is omitted.}
\end{figure}
